% !TITLE 晚登三山还望京邑
% !AUTHOR 谢朓
% !DYNASTY 南北朝
% !GENRE 五言古诗
% !ORDER 3

\begin{poem}
灞涘\textsuperscript{1}望长安，河阳\textsuperscript{2}视京县。
白日丽飞甍\textsuperscript{3}，参差皆可见。
余霞散成绮\textsuperscript{4}，澄江静如练\textsuperscript{5}。
喧鸟覆春洲\textsuperscript{6}，杂英满芳甸\textsuperscript{7}。
去矣方滞淫\textsuperscript{8}，怀哉罢欢宴。
佳期怅何许，泪下如流霰\textsuperscript{9}。
有情知望乡，谁能鬒\textsuperscript{10}不变？
\end{poem}

\begin{annotations}
\notetext{1}{灞涘：灞水之滨，在长安城外。这里是引用王粲《七哀诗》中“南登霸陵岸，回首望长安”的典故，表达对京城的回望与不舍。}
\notetext{2}{河阳：今河南孟州一带。诗人借用潘岳在河阳望洛阳的典故，将自己的离京之情与古人相印证。}
\notetext{3}{飞甍：高高翘起的屋脊。甍（méng），屋脊。阳光照在京城的飞檐上，高低错落的建筑在暮色中清晰可见。}
\notetext{4}{绮：有花纹的丝织品。天边的晚霞散开，像一匹绚烂的锦缎。一个“散”字写了霞光铺展的动态。}
\notetext{5}{练：白色的丝绸。清澈的江水静静地流淌，像一匹白练。谢朓这句“澄江静如练”是千古名句，李白读后赞不绝口。}
\notetext{6}{春洲：春天的沙洲。喧闹的鸟儿落满了江中的小岛，写尽了春天的生机。}
\notetext{7}{芳甸：开满鲜花的郊野。甸（diàn），郊外的田野。五颜六色的野花开遍了原野。}
\notetext{8}{滞淫：逗留、停滞不前。明知道要走了，却舍不得离开，脚步迈不出去。}
\notetext{9}{霰：雪珠、雪粒。眼泪像雪粒一样纷纷落下——不是一滴两滴，而是止不住的泪。}
\notetext{10}{鬒：黑发。鬒（zhěn），头发乌黑。思乡之苦让人一夜白头，谁能保持黑发不变呢？}
\end{annotations}

\begin{appreciation}
谢朓（tiǎo），南朝齐的诗人，也是李白最服气的诗人之一。李白夸他"解道澄江静如练，令人长忆谢玄晖"，记了一辈子——这首诗里就有那一句。

三山在建康（今南京）城外。谢朓要离京去外地做官，傍晚登三山回望京城，写下这首诗。

"灞涘望长安，河阳视京县"——开头借古人的姿势。王粲写"南登霸陵岸，回首望长安"，潘岳在河阳遥望京城。他说，我此刻的回望和古人一样。回望的动作一摆出来，全诗的调子就定了：舍不得走。

"白日丽飞甍"：夕阳明晃晃地照在京城翘起的屋脊上，高低错落都看得见。接着是那句千古名句——"余霞散成绮，澄江静如练"：晚霞在天边铺展，像一匹花纹绚烂的锦缎；江水安安静静地躺着，像一匹雪白的绸子。一个"散"字写霞光的流动，一个"静"字写江水的安详，一暖一冷，一天一地。十个字，把黄昏写尽了。他写景这么好，不是为了写景——他是要把京城的模样刻进眼里，马上看不到了。

"喧鸟覆春洲，杂英满芳甸"：鸟扑棱棱落满春天的沙洲，野花开遍郊野。越热闹，越衬出下一句"去矣方滞淫"——明明要走，脚步却钉在原地。"佳期怅何许，泪下如流霰"：再见遥遥无期，眼泪像雪珠落下。最后一句点题："有情知望乡，谁能鬒不变？"有情人都知道望乡的滋味，谁能不白头？鬒就是黑发。

还记得《行行重行行》吗？"思君令人老，岁月忽已晚"，是家里的人等游子等瘦了；谢朓是游子自己，站在山上看故乡看到白头。一个在家等，一个在路上望——等和望，原来都催人老。从古诗十九首到唐诗成片的望乡诗，这个主题中国人写了两千年。

妹妹，这首诗是给你以后出远门备着的。哪天你在外地想家了，翻出来读一遍，诗里的人会替你把话说出来。只是别学他望到白头——想家就回，机票钱哥哥出。
\end{appreciation}
